% LaTeX Strategic Analysis for AMC12B 2016 Problem 15
% Using the AMC12/COMC Problem Solving Framework

\documentclass[]{article}
\usepackage[utf8]{inputenc}
\usepackage{amsmath, amssymb, amsthm}
\usepackage{geometry}
\usepackage{xcolor}
\usepackage[most]{tcolorbox}
\usepackage{enumitem}
\usepackage{fancyhdr}

% Page setup
\geometry{margin=1in}
\pagestyle{fancy}
\fancyhf{}
\rhead{AMC12/COMC Problem Analysis}
\lhead{\today}

% Color definitions
\definecolor{problemconfig}{RGB}{230, 242, 255}    % Light blue
\definecolor{strategic}{RGB}{255, 230, 242}       % Light pink
\definecolor{tactical}{RGB}{242, 255, 230}        % Light green
\definecolor{techniques}{RGB}{255, 242, 230}      % Light yellow
\definecolor{verification}{RGB}{255, 230, 230}    % Light red
\definecolor{solution}{RGB}{230, 255, 255}        % Light cyan
\definecolor{competition}{RGB}{242, 242, 255}     % Light purple
\definecolor{gold}{RGB}{218, 165, 32}

% Box styles
\newtcolorbox{problemconfigbox}{
    colback=problemconfig,
    colframe=blue,
    title=Problem Configuration Analysis,
    fonttitle=\bfseries,
    arc=3pt,
    boxrule=1pt,
    breakable
}

\newtcolorbox{strategicbox}{
    colback=strategic,
    colframe=purple,
    title=Strategic Framework Application,
    fonttitle=\bfseries,
    arc=3pt,
    boxrule=1pt,
    breakable
}

\newtcolorbox{tacticalbox}{
    colback=tactical,
    colframe=green,
    title=Tactical Methodology Selection,
    fonttitle=\bfseries,
    arc=3pt,
    boxrule=1pt
}

\newtcolorbox{techniquesbox}{
    colback=techniques,
    colframe=orange,
    title=Conversion Techniques Application,
    fonttitle=\bfseries,
    arc=3pt,
    boxrule=1pt,
    breakable
}

\newtcolorbox{verificationbox}{
    colback=verification,
    colframe=red,
    title=Verification Strategy Design,
    fonttitle=\bfseries,
    arc=3pt,
    boxrule=1pt,
    breakable
}

\newtcolorbox{solutionbox}{
    colback=solution,
    colframe=cyan,
    title=Solution Roadmap,
    fonttitle=\bfseries,
    arc=3pt,
    boxrule=1pt,
    breakable
}

\newtcolorbox{competitionbox}{
    colback=competition,
    colframe=purple,
    title=Competition Strategy Integration,
    fonttitle=\bfseries,
    arc=3pt,
    boxrule=1pt,
    breakable
}

\newtcolorbox{keyinsightbox}{
    colback=yellow!20,
    colframe=gold,
    title=Key Insight,
    fonttitle=\bfseries,
    arc=3pt,
    boxrule=2pt,
    leftrule=5pt,
    breakable
}

\newtcolorbox{warningbox}{
    colback=red!10,
    colframe=red!50,
    title=Warning: Common Pitfall,
    fonttitle=\bfseries,
    arc=3pt,
    boxrule=2pt,
    leftrule=5pt,
    breakable
}

\newtcolorbox{tipbox}{
    colback=green!10,
    colframe=green!50,
    title=Pro Tip,
    fonttitle=\bfseries,
    arc=3pt,
    boxrule=2pt,
    leftrule=5pt,
    breakable
}

\begin{document}

\title{AMC12B 2016 Problem 15: Strategic Analysis}
\author{Competition Math Framework Analysis}
\date{\today}
\maketitle

\section*{Problem Statement}
All the numbers $2, 3, 4, 5, 6, 7$ are assigned to the six faces of a cube, one number to each face. For each of the eight vertices of the cube, a product of three numbers is computed, where the three numbers are the numbers assigned to the three faces that include that vertex. What is the greatest possible value of the sum of these eight products?

\textbf{Answer Choices:}
$\textbf{(A)}\ 312 \qquad \textbf{(B)}\ 343 \qquad \textbf{(C)}\ 625 \qquad \textbf{(D)}\ 729 \qquad \textbf{(E)}\ 1680$

\begin{problemconfigbox}
\subsection*{A. Problem Classification}
\begin{itemize}[leftmargin=*]
    \item \textbf{Problem Type}: To Find - determine maximum value
    \item \textbf{Difficulty Band}: Mid-band (Problem 15)
    \item \textbf{Competition Context}: AMC12 (75 min, 25 problems, 3 min target)
    \item \textbf{Mathematical Domain}: Combinatorics/Optimization with Algebraic Manipulation
\end{itemize}

\subsection*{B. Problem Structure Identification}
\begin{itemize}[leftmargin=*]
    \item \textbf{Given Information}: 
    \begin{itemize}
        \item Numbers 2, 3, 4, 5, 6, 7 assigned to cube faces (one per face)
        \item Cube has 8 vertices
        \item Each vertex touches exactly 3 faces
        \item Sum computed over all 8 vertex products
    \end{itemize}
    \item \textbf{Constraints}: 
    \begin{itemize}
        \item Each number used exactly once
        \item Assignment determines all vertex products
        \item Cube geometry: opposite faces never meet at same vertex
    \end{itemize}
    \item \textbf{Objective}: Maximize $\sum_{v=1}^{8} (\text{product at vertex } v)$
    \item \textbf{Hidden Assumptions}: 
    \begin{itemize}
        \item Cube has 3 pairs of opposite faces
        \item Each number appears in exactly 4 vertex products
        \item Symmetry suggests algebraic structure
    \end{itemize}
    \item \textbf{Answer Choice Leverage}: 
    \begin{itemize}
        \item (E) 1680 = $7 \times 6 \times 5 \times 8$ suspiciously large - likely trap
        \item (D) 729 = $9^3 = 3^6$ has interesting structure
        \item (B) 343 = $7^3$ also has power structure
        \item Can use choices to guide algebraic simplification
    \end{itemize}
\end{itemize}

\subsection*{C. Strategic Orientation}
\begin{itemize}[leftmargin=*]
    \item \textbf{Hypothesis}: Assignment of numbers to faces determines sum
    \item \textbf{Conclusion}: Find optimal assignment yielding maximum sum
    \item \textbf{Notation}: Let opposite face pairs be $(a,f)$, $(b,e)$, $(c,d)$
    \item \textbf{Intuitive Approach}: 
    \begin{itemize}
        \item This feels like an algebraic factorization problem
        \item Cube symmetry suggests the sum has clean form
        \item Try to express sum in terms of face values
    \end{itemize}
    \item \textbf{Cross-Domain Potential}: 
    \begin{itemize}
        \item Pure combinatorial search (inefficient)
        \item Algebraic manipulation via symmetry
        \item Optimization via AM-GM or calculus
    \end{itemize}
\end{itemize}
\end{problemconfigbox}

\begin{strategicbox}
\subsection*{A. Psychological Strategies}
\begin{itemize}[leftmargin=*]
    \item \textbf{Mental Toughness}: Problem appears complex initially, but algebraic structure likely simplifies dramatically
    \item \textbf{Creativity Cultivation}: Look for symmetric patterns rather than brute-forcing arrangements
    \item \textbf{Iterative Refinement}: If algebraic approach stalls, try small cases (e.g., numbers 1,2,3,4,5,6)
\end{itemize}

\subsection*{B. Investigation Strategies}
\begin{itemize}[leftmargin=*]
    \item \textbf{Startup Orientation}: Mid-band optimization problem with geometric constraint
    \item \textbf{Get Your Hands Dirty}: 
    \begin{itemize}
        \item Label cube faces and write explicit sum formula
        \item Use cube geometry: if face $a$ is opposite to $f$, they never share a vertex
    \end{itemize}
    \item \textbf{Penultimate Step}: Once sum is in factored form $(a+f)(b+e)(c+d)$, apply AM-GM or check answer choices
    \item \textbf{Wishful Thinking}: What if the sum factors into products of sums? Test this hypothesis!
    \item \textbf{Make It Easier}: Imagine numbers were 1,2,3,4,5,6 - would pattern be clearer?
    \item \textbf{Conjecture via Small Cases}: 
    \begin{itemize}
        \item For 2×2×2 scenario with 4 faces: would sum factor?
        \item Verify factorization works for simple assignment
    \end{itemize}
    \item \textbf{Visualization}: Draw cube, label faces, identify which triplets form at each vertex
\end{itemize}

\subsection*{C. Late-Band Strategic Playbook}
\begin{itemize}[leftmargin=*]
    \item \textbf{Answer-Choice Leverage}: 
    \begin{itemize}
        \item Test if $(a+f)(b+e)(c+d) = 729 = 9^3$ works
        \item This requires $(a+f) = (b+e) = (c+d) = 9$
        \item Numbers available: 2,3,4,5,6,7 with sum = 27
        \item Pairing: (2,7), (3,6), (4,5) gives $9 \times 9 \times 9 = 729$ $\checkmark$
    \end{itemize}
    \item \textbf{Penultimate Step}: Once factorization known, optimize $(a+f)(b+e)(c+d)$ subject to $a+b+c+d+e+f=27$
\end{itemize}

\subsection*{D. Argument Construction Strategy}
\begin{itemize}[leftmargin=*]
    \item \textbf{Direct Proof}: 
    \begin{enumerate}
        \item Express sum algebraically using cube structure
        \item Factor into product of sums
        \item Apply AM-GM: $\frac{(a+f)+(b+e)+(c+d)}{3}\geq\sqrt[3]{(a+f)(b+e)(c+d)}$
        \item Equality when $(a+f) = (b+e) = (c+d)$
    \end{enumerate}
\end{itemize}
\end{strategicbox}

\begin{tacticalbox}
\subsection*{A. Core Mathematical Tactics}
\begin{itemize}[leftmargin=*]
    \item \textbf{Symmetry Exploitation}:
    \begin{itemize}
        \item Cube has rotational symmetry
        \item Opposite faces play symmetric role
        \item Sum structure should reflect this symmetry
    \end{itemize}
    \item \textbf{Algebraic Manipulation}:
    \begin{itemize}
        \item Expand sum of 8 vertex products
        \item Look for common factors
        \item Factor completely using cube geometry
    \end{itemize}
    \item \textbf{Extreme Principle}:
    \begin{itemize}
        \item Maximum occurs at boundary of feasible region
        \item For optimization, check when sums are equal
    \end{itemize}
    \item \textbf{Inequality Application (AM-GM)}:
    \begin{itemize}
        \item For fixed sum $S$, product $xyz$ maximized when $x=y=z=S/3$
        \item Here: $(a+f)+(b+e)+(c+d) = 27$ fixed
        \item Maximum when $(a+f) = (b+e) = (c+d) = 9$
    \end{itemize}
\end{itemize}

\subsection*{B. Domain-Specific Tactics}
\begin{itemize}[leftmargin=*]
    \item \textbf{Combinatorics}: Each number appears in exactly 4 vertices
    \item \textbf{Optimization}: Convert discrete problem to continuous via AM-GM
    \item \textbf{Geometry}: Cube structure imposes constraint on which triplets can occur
\end{itemize}

\subsection*{C. Crossover Tactics}
\begin{itemize}[leftmargin=*]
    \item \textbf{Algebraic Geometry Connection}: Cube topology determines algebraic structure
    \item \textbf{Discrete to Continuous}: Use AM-GM from continuous optimization on discrete problem
\end{itemize}
\end{tacticalbox}

\begin{techniquesbox}
\subsection*{A. High-Probability Techniques (Recognition Index)}
\begin{itemize}[leftmargin=*]
    \item \textbf{Technique 1: Algebraic Factorization} [Confidence: 95\%]
    \begin{itemize}
        \item \textit{Recognition cue}: Sum over symmetric structure (cube vertices)
        \item \textit{Application}: Label opposite faces $(a,f)$, $(b,e)$, $(c,d)$
        \item \textit{Result}: Sum = $(a+f)(b+e)(c+d)$
        \item \textit{Why it works}: Cube geometry ensures each opposite pair appears together 4 times
    \end{itemize}
    
    \item \textbf{Technique 2: AM-GM Inequality} [Confidence: 90\%]
    \begin{itemize}
        \item \textit{Recognition cue}: Optimization problem with product structure
        \item \textit{Application}: $\frac{(a+f)+(b+e)+(c+d)}{3} \geq \sqrt[3]{(a+f)(b+e)(c+d)}$
        \item \textit{Result}: Maximum when sums are equal
        \item \textit{Why it works}: Product of fixed sum maximized at equality
    \end{itemize}
    
    \item \textbf{Technique 3: Optimization via Pairing} [Confidence: 85\%]
    \begin{itemize}
        \item \textit{Recognition cue}: Need to pair elements to maximize product
        \item \textit{Application}: Pair smallest with largest: (2,7), (3,6), (4,5)
        \item \textit{Result}: Equal sums of 9 each
        \item \textit{Why it works}: Balancing creates equal factors
    \end{itemize}
\end{itemize}

\subsection*{B. Technique Integration Strategy}
\begin{enumerate}
    \item \textbf{Phase 1: Algebraic Setup} (2 min)
    \begin{itemize}
        \item Express sum in terms of face labels
        \item Identify factorization pattern
    \end{itemize}
    
    \item \textbf{Phase 2: Optimization} (1 min)
    \begin{itemize}
        \item Apply AM-GM to establish maximum condition
        \item Determine optimal pairing strategy
    \end{itemize}
    
    \item \textbf{Phase 3: Verification} (1 min)
    \begin{itemize}
        \item Check pairing achieves equal sums
        \item Compute final answer: $9^3 = 729$
    \end{itemize}
\end{enumerate}

\subsection*{C. Technique Conflict Resolution}
\begin{itemize}[leftmargin=*]
    \item \textbf{Potential Conflict}: Brute-force enumeration vs. algebraic approach
    \begin{itemize}
        \item \textit{Resolution}: Algebraic approach much faster (4 min vs. 15+ min)
        \item \textit{Backup plan}: If factorization unclear, try answer-choice testing
    \end{itemize}
\end{itemize}
\end{techniquesbox}

\begin{verificationbox}
\subsection*{A. Verification Level Selection}
\begin{itemize}[leftmargin=*]
    \item \textbf{Selected Level}: Level 4 (Comprehensive) - mid-band problem worth full verification
    \item \textbf{Decision Rationale}:
    \begin{itemize}
        \item Mid-band position (P15) - should be solved correctly
        \item Clean algebraic structure reduces error risk
        \item Answer value 729 = $9^3$ is memorizable constant
    \end{itemize}
\end{itemize}

\subsection*{B. Primary Verification Methods}
\begin{enumerate}
    \item \textbf{Algebraic Verification}: 
    \begin{itemize}
        \item Verify factorization: Expand $(a+f)(b+e)(c+d)$ and confirm it equals the sum of 8 vertex products
        \item For assignment (2,7), (3,6), (4,5): compute $9 \times 9 \times 9 = 729$ $\checkmark$
        \item Check one vertex explicitly: With $a=7, b=6, c=5$ we get vertex $abc = 7 \times 6 \times 5 = 210$
    \end{itemize}
    
    \item \textbf{Bound Checking}:
    \begin{itemize}
        \item Upper bound: $(7+2)(6+3)(5+4) = 9 \times 9 \times 9 = 729$ $\checkmark$
        \item Lower bound: Unequal pairings give smaller products
        \item Example: $(7+3)(6+2)(5+4) = 10 \times 8 \times 9 = 720 < 729$ $\checkmark$
        \item Another: $(7+4)(6+5)(3+2) = 11 \times 11 \times 5 = 605 < 729$ $\checkmark$
    \end{itemize}
    
    \item \textbf{Answer Choice Consistency}:
    \begin{itemize}
        \item 729 = $9^3$ appears as choice (D)
        \item 1680 = $7 \times 6 \times 5 \times 8$ represents wrong model (all max at one vertex)
        \item Our answer aligns with structural expectation
    \end{itemize}
    
    \item \textbf{Small Case Verification}:
    \begin{itemize}
        \item Test with numbers 1,2,3,4,5,6 (sum = 21)
        \item Optimal pairing: (1,6), (2,5), (3,4) gives $7 \times 7 \times 7 = 343$
        \item Note: 343 = $7^3$ appears as choice (B) - this confirms our pattern! $\checkmark$
    \end{itemize}
\end{enumerate}

\subsection*{C. Confidence Calibration}
\begin{itemize}[leftmargin=*]
    \item \textbf{Initial Confidence}: 85\% after finding factorization
    \item \textbf{Post-Verification Confidence}: 98\%
    \item \textbf{Remaining Uncertainty}: Possible algebraic expansion error (2\%)
    \item \textbf{Flag for Review?}: No - high confidence with multiple verification methods
\end{itemize}
\end{verificationbox}

\begin{solutionbox}
\subsection*{Step-by-Step Solution Plan}

\textbf{Step 1: Understand Cube Geometry} [30 seconds]
\begin{itemize}
    \item Cube has 6 faces, 8 vertices, 12 edges
    \item Each vertex is formed by 3 faces meeting
    \item Opposite faces never meet at the same vertex
    \item \textit{Checkpoint}: Can you identify the 3 pairs of opposite faces?
\end{itemize}

\textbf{Step 2: Label Faces Algebraically} [30 seconds]
\begin{itemize}
    \item Let opposite face pairs be: $(a,f)$, $(b,e)$, $(c,d)$
    \item These 6 variables represent numbers 2,3,4,5,6,7 in some order
    \item \textit{Checkpoint}: Do you see why opposite faces can't share a vertex?
\end{itemize}

\textbf{Step 3: Express Sum of Vertex Products} [1 minute]
\begin{itemize}
    \item List all 8 vertices by the 3 faces that meet there:
    \begin{itemize}
        \item $abc, abd, aec, aed$ (four vertices involving face $a$)
        \item $fbc, fbd, fec, fed$ (four vertices involving face $f$)
    \end{itemize}
    \item Sum = $abc + abd + aec + aed + fbc + fbd + fec + fed$
    \item \textit{Key observation}: Notice opposite faces $b$ and $e$ never appear together, same for $c$ and $d$
    \item \textit{Checkpoint}: Do you have all 8 vertices with correct face combinations?
\end{itemize}

\textbf{Step 4: Factor the Sum} [1 minute]
\begin{itemize}
    \item Group by common factors involving $a$ and $f$:
    \begin{align*}
        \text{Sum} &= abc + abd + aec + aed + fbc + fbd + fec + fed \\
        &= a(bc + bd + ec + ed) + f(bc + bd + ec + ed) \\
        &= (a+f)(bc + bd + ec + ed)
    \end{align*}
    \item Now factor the expression $(bc + bd + ec + ed)$:
    \begin{align*}
        bc + bd + ec + ed &= b(c+d) + e(c+d) \\
        &= (b+e)(c+d)
    \end{align*}
    \item Therefore: $\boxed{\text{Sum} = (a+f)(b+e)(c+d)}$
    \item \textit{Critical insight}: Sum factors into product of opposite face sums!
    \item \textit{Verification}: Expand $(a+f)(b+e)(c+d)$ to confirm it gives all 8 vertex products
    \item \textit{Checkpoint}: Each opposite face pair (like $a$ and $f$) appears together in exactly 4 vertices
\end{itemize}

\textbf{Step 5: Apply AM-GM Inequality} [45 seconds]
\begin{itemize}
    \item We know: $a+b+c+d+e+f = 2+3+4+5+6+7 = 27$
    \item Therefore: $(a+f)+(b+e)+(c+d) = 27$
    \item By AM-GM: $\frac{(a+f)+(b+e)+(c+d)}{3} \geq \sqrt[3]{(a+f)(b+e)(c+d)}$
    \item So: $9 \geq \sqrt[3]{(a+f)(b+e)(c+d)}$
    \item Thus: $(a+f)(b+e)(c+d) \leq 729$
    \item Equality when $(a+f) = (b+e) = (c+d) = 9$
\end{itemize}

\textbf{Step 6: Find Optimal Pairing} [30 seconds]
\begin{itemize}
    \item Need three pairs from $\{2,3,4,5,6,7\}$ each summing to 9:
    \begin{itemize}
        \item $2+7 = 9$ $\checkmark$
        \item $3+6 = 9$ $\checkmark$
        \item $4+5 = 9$ $\checkmark$
    \end{itemize}
    \item This pairing achieves the maximum!
\end{itemize}

\textbf{Step 7: Compute Final Answer} [15 seconds]
\begin{itemize}
    \item Maximum sum = $9 \times 9 \times 9 = 729$
    \item This is answer choice (D)
\end{itemize}

\subsection*{Alternative Approaches}

\textbf{Alternative 1: Direct Computation} [Time: 8-10 min]
\begin{itemize}
    \item Try specific assignment: opposite faces (2,7), (3,6), (4,5)
    \item Compute all 8 vertex products explicitly
    \item Sum: $7 \cdot 6 \cdot 5 + 7 \cdot 6 \cdot 4 + 7 \cdot 3 \cdot 5 + 7 \cdot 3 \cdot 4$\\
          $+ 2 \cdot 6 \cdot 5 + 2 \cdot 6 \cdot 4 + 2 \cdot 3 \cdot 5 + 2 \cdot 3 \cdot 4$
    \item Calculate: $210+168+105+84+60+48+30+24 = 729$
    \item \textit{Downside}: Time-consuming and error-prone
\end{itemize}

\textbf{Alternative 2: Answer Choice Testing} [Time: 5-7 min]
\begin{itemize}
    \item Test (E) 1680: Impossible - exceeds maximum possible sum
    \item Test (D) 729 = $9^3$: Suggests equal factors of 9
    \item Verify pairing (2,7), (3,6), (4,5) achieves this
    \item Can compute directly: $9 \times 9 \times 9 = 729$
    \item \textit{Advantage}: Works well if algebraic approach unclear
\end{itemize}

\subsection*{Common Pitfalls}

\begin{enumerate}
    \item \textbf{Pitfall}: Trying to maximize individual vertex products rather than sum
    \begin{itemize}
        \item \textit{Why wrong}: Can't have $7 \times 6 \times 5$ at all vertices
        \item \textit{Trap}: Answer choice (E) 1680 = $8 \times 210$ exploits this error
    \end{itemize}
    
    \item \textbf{Pitfall}: Not recognizing factorization pattern
    \begin{itemize}
        \item \textit{Consequence}: Leads to brute-force enumeration
        \item \textit{Recovery}: Use answer choice (D) to guide pairing strategy
    \end{itemize}
    
    \item \textbf{Pitfall}: Forgetting cube has opposite faces
    \begin{itemize}
        \item \textit{Error}: Assuming any 3 faces can meet at vertex
        \item \textit{Check}: Verify geometric constraint in setup
    \end{itemize}
    
    \item \textbf{Pitfall}: Arithmetic errors in factorization
    \begin{itemize}
        \item \textit{Prevention}: Verify factorization with small example
        \item \textit{Check}: Expand $(a+f)(b+d)(c+e)$ for one specific case
    \end{itemize}
\end{enumerate}

\subsection*{Time Management}
\begin{itemize}
    \item \textbf{Target Time}: 4 minutes (appropriate for P15)
    \item \textbf{Actual Breakdown}:
    \begin{itemize}
        \item Setup \& labeling: 1 min
        \item Algebraic factorization: 1.5 min
        \item AM-GM application: 0.75 min
        \item Verification: 0.75 min
        \item \textbf{Total}: ~4 minutes $\checkmark$
    \end{itemize}
    \item \textbf{Abandon Threshold}: If no progress after 6 minutes, guess and move on
\end{itemize}
\end{solutionbox}

\begin{competitionbox}
\subsection*{A. Time Allocation Strategy}
\begin{itemize}[leftmargin=*]
    \item \textbf{Problem Position}: P15 (mid-band, borderline to late)
    \item \textbf{Target Time}: 4 minutes
    \item \textbf{Maximum Time}: 6 minutes before abandoning
    \item \textbf{Time Checkpoints}:
    \begin{itemize}
        \item 2 min: Should have factorization insight
        \item 4 min: Should have answer and be verifying
        \item 6 min: Move on if still stuck
    \end{itemize}
\end{itemize}

\subsection*{B. Problem Prioritization Strategy}
\begin{itemize}[leftmargin=*]
    \item \textbf{When to Attempt}: 
    \begin{itemize}
        \item After completing P1-P12 (easier problems)
        \item If comfortable with optimization and factorization
        \item If you recognize AM-GM pattern quickly
    \end{itemize}
    \item \textbf{Skip Indicators}:
    \begin{itemize}
        \item Uncomfortable with 3D geometry
        \item Algebraic factorization not a strength
        \item Running low on time
    \end{itemize}
\end{itemize}

\subsection*{C. Risk Management}
\begin{itemize}[leftmargin=*]
    \item \textbf{Confidence Threshold}: Medium-high (70-85\%)
    \begin{itemize}
        \item Algebraic structure is clean
        \item AM-GM is standard technique
        \item Answer 729 = $9^3$ is memorable
    \end{itemize}
    \item \textbf{Soft Abandon Triggers}:
    \begin{itemize}
        \item Can't see factorization after 3 minutes
        \item Algebraic manipulation becoming messy
        \item Made calculation error twice
    \end{itemize}
    \item \textbf{Hard Abandon Triggers}:
    \begin{itemize}
        \item No progress after 6 minutes
        \item Emotional frustration setting in
        \item Other easier problems remain
    \end{itemize}
\end{itemize}

\subsection*{D. Optimization Strategies}
\begin{itemize}[leftmargin=*]
    \item \textbf{Speed Techniques}:
    \begin{itemize}
        \item Recognize $729 = 9^3$ immediately from answer choices
        \item This hints at factorization $(a+f)(b+e)(c+d)$ with equal sums
        \item Skip full AM-GM derivation; just verify pairing works
    \end{itemize}
    \item \textbf{Verification Shortcuts}:
    \begin{itemize}
        \item Don't need to expand full product
        \item Just check one alternative pairing gives smaller value
        \item Example: $(10)(8)(9) = 720 < 729$ $\checkmark$
    \end{itemize}
\end{itemize}
\end{competitionbox}

\begin{keyinsightbox}
\textbf{Key Insight}: The sum of the 8 vertex products factors beautifully as $(a+f)(b+e)(c+d)$ where $(a,f)$, $(b,e)$, $(c,d)$ are opposite face pairs. This factorization transforms a complex combinatorial optimization into a simple AM-GM inequality problem. The maximum occurs when opposite faces sum equally: $2+7=3+6=4+5=9$, giving $9^3 = 729$.
\end{keyinsightbox}

\begin{warningbox}
\textbf{Warning: Common Pitfall}

The trap answer (E) 1680 = $7 \times 6 \times 5 \times 8$ exploits the misconception that you can maximize each vertex independently. Students mistakenly think: "Put 7,6,5 at one vertex to get $7 \times 6 \times 5 = 210$, then do this at all 8 vertices." This is IMPOSSIBLE because each face touches 4 vertices, not 8. The cube's geometry constrains which triplets can occur simultaneously. Always verify your answer respects the problem's structural constraints!
\end{warningbox}

\begin{tipbox}
\textbf{Pro Tip: Competition Time-Saver}

When you see an optimization problem on a cube with product structure, immediately think "opposite faces" and "factorization." The answer choices are your friend: $729 = 9^3$ screams "three equal factors of 9." Work backwards: if you need three pairs summing to 9 from $\{2,3,4,5,6,7\}$, there's only one way: (2,7), (3,6), (4,5). Verify this works, bubble (D), and move on in under 3 minutes!
\end{tipbox}

\section*{Final Answer}
\boxed{\textbf{(D)}\ 729}

\section*{Confidence Assessment}
\begin{itemize}[leftmargin=*]
    \item \textbf{Confidence Level}: 98\% - Multiple verification methods confirm answer
    \item \textbf{Verification Applied}: Level 4 (Comprehensive)
    \begin{itemize}
        \item Algebraic factorization verified
        \item AM-GM optimality confirmed
        \item Alternative pairing tested (gave 720 < 729)
        \item Small case (numbers 1-6) pattern confirmed
    \end{itemize}
    \item \textbf{Flag for Review}: No - extremely high confidence
    \item \textbf{Time Spent}: 4 minutes (matches target for P15)
\end{itemize}

\section*{Learning Outcomes}
\begin{enumerate}
    \item \textbf{Technique Mastery}: Factorization of symmetric sums over cube vertices
    \item \textbf{Pattern Recognition}: $n^3$ in answer choices suggests equal factors
    \item \textbf{Optimization Strategy}: AM-GM for products with fixed sum constraint
    \item \textbf{Geometric Intuition}: Opposite faces structure determines algebraic form
    \item \textbf{Verification Skill}: Multiple independent methods increase confidence
\end{enumerate}

\end{document}